\chapter{Der Weg zur Praxisprojektstelle}\label{ch:weg_zur_pp}
Im August 2024 habe ich mich bei der lise GmbH als Werkstudent in der Software Entwicklung beworben, unter der Vorraussetzung dass ich dort
mein Praxisprojekt absolvieren kann, da dies in naher Zukunft auf mich zu kommen würde.\\
Von Vorteil war hier meine vorherige Erfahrung als Software Entwickler bei meinen vorherigen Arbeitgebern wodurch wir ein
Thema finden konnten dass der lise GmbH sowohl Mehrwert bring als auch in den Rahmen des Praxisprojekts passt.\\
Das ist nicht ganz leicht gewesen da der übliche Projektalltag von den Vorgaben der Hochschule, ein abgestecktes Projekt zu behandeln, abweicht.
Trotzdessen konnte sich ein passendes Projekt finden lassen mit dem alle Parteien zufrieden sind.

\chapter{Unternehmensbeschreibung}\label{ch:unternehmen}
\section{Unternehmensdaten}\label{sec:nternehmensdaten}
lise GmbH

\textbf{Hauptsitz:} Rudi-Conin-Straße 5, 50829 Köln\\
\textbf{Gründungsjahr:} 1999\\
\textbf{Mitarbeiter: } 100\\
\textbf{Registernummer: } HRB 61293\\
\textbf{Bilanzsumme 2023: } 4.907.296,24 EUR \autocite{unternehmensregister2024}\\
\textbf{Website:} \url{https://www.lise.de/}\\
\textbf{Branchenzugehörigkeit:} Softwareentwicklungs- und IT-Consultings-Dienstleistungen\\
\textbf{Logo:}
\begin{figure}[H]
    \begin{center}
        \includegraphics[width=3cm]{bilder/lise_logo_web.png}
        \caption{Logo lise GmbH}\label{fig:lise}
    \end{center}
\end{figure}

\section{Dienstleistungen}\label{sec:dienstleistungen}
Die lise GmbH bietet maßgeschneiderte Softwarelösungen für die individuellen Ansprüche ihrer Kunden.\\
Diese erstrecken sich neben der Softwareentwicklung auch über Künstliche Intelligenz, UX-Design, Scrum,
Softwaremodernisierung, Sharepoint DevOps und generelle Beratung betüglich IT-Strategie.\autocite[/ueber-uns/unternehmen]{LiseWeb}

Die Kunden der lise sind aus diversen Industrien z.B Großhandel, Lebensmittel, Behörden und natürlich auch der IT-Branche.




% \section{Zitate}\label{sec:zitate}
% Nach~\cite{Bartels2008} wird ein Beispieltext vor dem ersten Zitat in einem Kapitel benötigt.\\
% \textcite{Yoshida1999} sagen hier steht ein Text. \\
% \textcite*{Burns2003} sagen hier steht ein Text. \\
% \enquote{Hier steht ein Text.}\autocite{Bartels2008} \\
% \enquote{Hier steht ein Text.}\autocite[Vgl.][]{Bartels2008} \\
% \enquote{Hier steht ein Text.}\autocite[][S. 200]{Bruno2009} \\
% \enquote{Hier steht ein Text.}\autocite*[][S. 200]{Bruno2009} \\
% \enquote{Hier steht ein Text.}\autocite{Wikipedia2011, Bruno2009}

% \subsection{Abkürzungen}\label{subsec:abkuerzungen}
% \gls{but}\\
% \gls{cdma}\\
% \gls{gsm}\\
% \gls{ic}\\
% \gls{lh2}\\
% \gls{lox}\\
% \gls{na}\\
% \gls{nad+}\\
% \gls{nua}\\
% \gls{tdma}\\
% \gls{ua}

% \subsection{Quellcode}\label{subsec:quellcode}
% \begin{figure}[H]
%     \begin{center}
%         \begin{minted}{c}
% int main() {
%   printf("hello, world");
%   return 0;
% }
%         \end{minted}
%         \caption{Hier steht eingebetteter Quellcode}\label{fig:hier-steht-eingebetteter-quellcode}
%     \end{center}
% \end{figure}

% \begin{figure}[H]
%     \begin{center}
%         \inputminted{python}{anhang/example.py}
%         \caption{Dieser Quellcode ist in einer externen Datei ausgelagert}\label{fig:dieser-quellcode-ist-in-einer-externen-datei-ausgelagert}
%     \end{center}
% \end{figure}
